
\documentclass{article}
\usepackage{colortbl}
\usepackage{makecell}
\usepackage{multirow}
\usepackage{supertabular}

\begin{document}

\newcounter{utterance}

\centering \large Interaction Transcript for game `cladder', experiment `full\_v1.5\_default', episode 2979 with Qwen3.5{-}9B.
\vspace{24pt}

{ \footnotesize  \setcounter{utterance}{1}
\setlength{\tabcolsep}{0pt}
\begin{supertabular}{c@{$\;$}|p{.15\linewidth}@{}p{.15\linewidth}p{.15\linewidth}p{.15\linewidth}p{.15\linewidth}p{.15\linewidth}}
    \# & \multicolumn{2}{c}{Player} && \multicolumn{2}{c}{Game Master} \\
    \hline

    \theutterance \stepcounter{utterance}  
    & & & \multicolumn{4}{p{0.6\linewidth}}{
        \cellcolor[rgb]{0.9,0.9,0.9}{
            \makecell[{{p{\linewidth}}}]{
                \texttt{\tiny{[P1$\langle$GM]}}
                \texttt{You are an expert in causal inference. The following question is not a typical commonsense query, but rather a meticulously designed question created by a professor specializing in causal inference, intended to assess the students' mastery of the course content.} \\
\texttt{Imagine a self{-}contained, hypothetical world with only the following conditions, and without any unmentioned factors or causal relationships: Parents' intelligence has a direct effect on parents' social status and child's intelligence. Other unobserved factors has a direct effect on parents' social status and child's intelligence. Parents' social status has a direct effect on child's intelligence. Other unobserved factors is unobserved. We know that intelligent parents and confounder active causes high parental social status. intelligent parents and confounder active and high parental social status causes intelligent child. We observed confounder active. Would the child is intelligent if intelligent parents instead of unintelligent parents?} \\
\texttt{Start your answer with 'Yes' or 'No', followed by additional reasoning or evidence to support your explanation.} \\
            }
        }
    }
     \\ \\

    \theutterance \stepcounter{utterance}  
    & \multicolumn{4}{p{0.6\linewidth}}{
        \cellcolor[rgb]{0.561,0.737,0.561}{
            \makecell[{{p{\linewidth}}}]{
                \texttt{\tiny{[P1$\rangle$GM]}}
                \texttt{No} \\
\\ 
\texttt{**Reasoning:**} \\
\\ 
\texttt{In this hypothetical world, the causal structure can be represented as follows:} \\
\texttt{1.  **Parents' Intelligence (\$P\_I\$)** \$\textbackslash{}rightarrow\$ **Parents' Social Status (\$P\_\{SS\}\$)**} \\
\texttt{2.  **Parents' Intelligence (\$P\_I\$)** \$\textbackslash{}rightarrow\$ **Child's Intelligence (\$C\_I\$)**} \\
\texttt{3.  **Confounder Active (\$C\$)** \$\textbackslash{}rightarrow\$ **Parents' Social Status (\$P\_\{SS\}\$)**} \\
\texttt{4.  **Confounder Active (\$C\$)** \$\textbackslash{}rightarrow\$ **Child's Intelligence (\$C\_I\$)**} \\
\texttt{5.  **Parents' Social Status (\$P\_\{SS\}\$)** \$\textbackslash{}rightarrow\$ **Child's Intelligence (\$C\_I\$)**} \\
\\ 
\texttt{The question asks for the counterfactual: Would the child be intelligent if the parents were intelligent, given that we observed the confounder active (\$C=1\$) and presumably the parents were originally unintelligent (\$P\_I=0\$)?} \\
\\ 
\texttt{To answer this, we must determine the value of \$C\_I\$ in the world where \$P\_I=1\$, while holding the confounder \$C\$ fixed at its observed value (\$C=1\$).} \\
\\ 
\texttt{The total effect of Parents' Intelligence on Child's Intelligence is mediated through two paths:} \\
\texttt{1.  **Direct Path:** \$P\_I \textbackslash{}rightarrow C\_I\$} \\
\texttt{2.  **Indirect Path:** \$P\_I \textbackslash{}rightarrow P\_\{SS\} \textbackslash{}rightarrow} \\
            }
        }
    }
    & & \\ \\

    \theutterance \stepcounter{utterance}  
    & & & \multicolumn{2}{p{0.3\linewidth}}{
        \cellcolor[rgb]{0.9,0.9,0.9}{
            \makecell[{{p{\linewidth}}}]{
                \texttt{\tiny{[GM$|$GM]}}
                \texttt{No} \\
            }
        }
    }
    & & \\ \\

    \theutterance \stepcounter{utterance}  
    & & & \multicolumn{2}{p{0.3\linewidth}}{
        \cellcolor[rgb]{0.9,0.9,0.9}{
            \makecell[{{p{\linewidth}}}]{
                \texttt{\tiny{[GM$|$GM]}}
                \texttt{yes} \\
            }
        }
    }
    & & \\ \\

    \theutterance \stepcounter{utterance}  
    & & & \multicolumn{2}{p{0.3\linewidth}}{
        \cellcolor[rgb]{0.9,0.9,0.9}{
            \makecell[{{p{\linewidth}}}]{
                \texttt{\tiny{[GM$|$GM]}}
                \texttt{game\_result = LOSE} \\
            }
        }
    }
    & & \\ \\

\end{supertabular}
}

\end{document}
